%% Generated by Sphinx.
\def\sphinxdocclass{report}
\documentclass[letterpaper,10pt,english]{sphinxmanual}
\ifdefined\pdfpxdimen
   \let\sphinxpxdimen\pdfpxdimen\else\newdimen\sphinxpxdimen
\fi \sphinxpxdimen=.75bp\relax
\ifdefined\pdfimageresolution
    \pdfimageresolution= \numexpr \dimexpr1in\relax/\sphinxpxdimen\relax
\fi
\newdimen\sphinxremdimen\sphinxremdimen = 10pt
%% let collapsible pdf bookmarks panel have high depth per default
\PassOptionsToPackage{bookmarksdepth=5}{hyperref}
%% turn off hyperref patch of \index as sphinx.xdy xindy module takes care of
%% suitable \hyperpage mark-up, working around hyperref-xindy incompatibility
\PassOptionsToPackage{hyperindex=false}{hyperref}
%% memoir class requires extra handling
\makeatletter\@ifclassloaded{memoir}
{\ifdefined\memhyperindexfalse\memhyperindexfalse\fi}{}\makeatother

\PassOptionsToPackage{booktabs}{sphinx}
\PassOptionsToPackage{colorrows}{sphinx}

\PassOptionsToPackage{warn}{textcomp}

\catcode`^^^^00a0\active\protected\def^^^^00a0{\leavevmode\nobreak\ }
\usepackage{cmap}
\usepackage{fontspec}
\defaultfontfeatures[\rmfamily,\sffamily,\ttfamily]{}
\usepackage{amsmath,amssymb,amstext}
\usepackage{polyglossia}
\setmainlanguage{english}



\setmainfont{FreeSerif}[
  Extension      = .otf,
  UprightFont    = *,
  ItalicFont     = *Italic,
  BoldFont       = *Bold,
  BoldItalicFont = *BoldItalic
]
\setsansfont{FreeSans}[
  Extension      = .otf,
  UprightFont    = *,
  ItalicFont     = *Oblique,
  BoldFont       = *Bold,
  BoldItalicFont = *BoldOblique,
]
\setmonofont{FreeMono}[Scale=0.9,
  Extension      = .otf,
  UprightFont    = *,
  ItalicFont     = *Oblique,
  BoldFont       = *Bold,
  BoldItalicFont = *BoldOblique,
]



\usepackage[Bjarne]{fncychap}
\usepackage{sphinx}

\fvset{fontsize=auto}
\usepackage{geometry}


% Include hyperref last.
\usepackage{hyperref}
% Fix anchor placement for figures with captions.
\usepackage{hypcap}% it must be loaded after hyperref.
% Set up styles of URL: it should be placed after hyperref.
\urlstyle{same}

\addto\captionsenglish{\renewcommand{\contentsname}{Contents:}}

\usepackage{sphinxmessages}
\setcounter{tocdepth}{1}



\title{ALU Documentation}
\date{Jun 19, 2026}
\release{1.0}
\author{Shlok Jha}
\newcommand{\sphinxlogo}{\vbox{}}
\renewcommand{\releasename}{Release}
\makeindex
\begin{document}

\pagestyle{empty}
\sphinxmaketitle
\pagestyle{plain}
\sphinxtableofcontents
\pagestyle{normal}
\phantomsection\label{\detokenize{index::doc}}


\sphinxstepscope


\chapter{Introduction}
\label{\detokenize{alu_intro:introduction}}\label{\detokenize{alu_intro::doc}}
\sphinxAtStartPar
The alu module is the Arithmetic Logic Unit (ALU)

\sphinxAtStartPar
It performs arithmetic, logical, shift and comparison
operations during instruction execution.

\sphinxAtStartPar
The ALU receives operands from the execution stage and
returns the processed result.

\sphinxstepscope


\chapter{Module Declaration}
\label{\detokenize{alu_moduledeclaration:module-declaration}}\label{\detokenize{alu_moduledeclaration::doc}}
\sphinxAtStartPar
The ALU module is declared as:

\begin{sphinxVerbatim}[commandchars=\\\{\}]
\PYG{n}{module} \PYG{n}{alu}
\end{sphinxVerbatim}

\sphinxstepscope


\chapter{Ports Description}
\label{\detokenize{alu_ports:ports-description}}\label{\detokenize{alu_ports::doc}}

\section{Inputs}
\label{\detokenize{alu_ports:inputs}}

\begin{savenotes}\sphinxattablestart
\sphinxthistablewithglobalstyle
\centering
\begin{tabulary}{\linewidth}[t]{TTT}
\sphinxtoprule
\begin{varwidth}[t]{\sphinxcolwidth{1}{3}}
\sphinxstyletheadfamily \sphinxAtStartPar
Signal
\sphinxbeforeendvarwidth
\end{varwidth}%
&\begin{varwidth}[t]{\sphinxcolwidth{1}{3}}
\sphinxstyletheadfamily \sphinxAtStartPar
Width
\sphinxbeforeendvarwidth
\end{varwidth}%
&\begin{varwidth}[t]{\sphinxcolwidth{1}{3}}
\sphinxstyletheadfamily \sphinxAtStartPar
Description
\sphinxbeforeendvarwidth
\end{varwidth}%
\\
\sphinxmidrule
\sphinxtableatstartofbodyhook\begin{varwidth}[t]{\sphinxcolwidth{1}{3}}
\sphinxAtStartPar
operator\_i
\sphinxbeforeendvarwidth
\end{varwidth}%
&\begin{varwidth}[t]{\sphinxcolwidth{1}{3}}
\sphinxAtStartPar
ALU\_OP\_WIDTH
\sphinxbeforeendvarwidth
\end{varwidth}%
&\begin{varwidth}[t]{\sphinxcolwidth{1}{3}}
\sphinxAtStartPar
Operation selector generated by decoder
\sphinxbeforeendvarwidth
\end{varwidth}%
\\
\sphinxhline\begin{varwidth}[t]{\sphinxcolwidth{1}{3}}
\sphinxAtStartPar
operand\_a\_i
\sphinxbeforeendvarwidth
\end{varwidth}%
&\begin{varwidth}[t]{\sphinxcolwidth{1}{3}}
\sphinxAtStartPar
32 bits
\sphinxbeforeendvarwidth
\end{varwidth}%
&\begin{varwidth}[t]{\sphinxcolwidth{1}{3}}
\sphinxAtStartPar
First operand
\sphinxbeforeendvarwidth
\end{varwidth}%
\\
\sphinxhline\begin{varwidth}[t]{\sphinxcolwidth{1}{3}}
\sphinxAtStartPar
operand\_b\_i
\sphinxbeforeendvarwidth
\end{varwidth}%
&\begin{varwidth}[t]{\sphinxcolwidth{1}{3}}
\sphinxAtStartPar
32 bits
\sphinxbeforeendvarwidth
\end{varwidth}%
&\begin{varwidth}[t]{\sphinxcolwidth{1}{3}}
\sphinxAtStartPar
Second operand
\sphinxbeforeendvarwidth
\end{varwidth}%
\\
\sphinxbottomrule
\end{tabulary}
\sphinxtableafterendhook\par
\sphinxattableend\end{savenotes}


\section{Outputs}
\label{\detokenize{alu_ports:outputs}}

\begin{savenotes}\sphinxattablestart
\sphinxthistablewithglobalstyle
\centering
\begin{tabulary}{\linewidth}[t]{TTT}
\sphinxtoprule
\begin{varwidth}[t]{\sphinxcolwidth{1}{3}}
\sphinxstyletheadfamily \sphinxAtStartPar
Signal
\sphinxbeforeendvarwidth
\end{varwidth}%
&\begin{varwidth}[t]{\sphinxcolwidth{1}{3}}
\sphinxstyletheadfamily \sphinxAtStartPar
Width
\sphinxbeforeendvarwidth
\end{varwidth}%
&\begin{varwidth}[t]{\sphinxcolwidth{1}{3}}
\sphinxstyletheadfamily \sphinxAtStartPar
Description
\sphinxbeforeendvarwidth
\end{varwidth}%
\\
\sphinxmidrule
\sphinxtableatstartofbodyhook\begin{varwidth}[t]{\sphinxcolwidth{1}{3}}
\sphinxAtStartPar
adder\_result\_o
\sphinxbeforeendvarwidth
\end{varwidth}%
&\begin{varwidth}[t]{\sphinxcolwidth{1}{3}}
\sphinxAtStartPar
32 bits
\sphinxbeforeendvarwidth
\end{varwidth}%
&\begin{varwidth}[t]{\sphinxcolwidth{1}{3}}
\sphinxAtStartPar
Result of internal adder
\sphinxbeforeendvarwidth
\end{varwidth}%
\\
\sphinxhline\begin{varwidth}[t]{\sphinxcolwidth{1}{3}}
\sphinxAtStartPar
result\_o
\sphinxbeforeendvarwidth
\end{varwidth}%
&\begin{varwidth}[t]{\sphinxcolwidth{1}{3}}
\sphinxAtStartPar
32 bits
\sphinxbeforeendvarwidth
\end{varwidth}%
&\begin{varwidth}[t]{\sphinxcolwidth{1}{3}}
\sphinxAtStartPar
Final ALU output
\sphinxbeforeendvarwidth
\end{varwidth}%
\\
\sphinxhline\begin{varwidth}[t]{\sphinxcolwidth{1}{3}}
\sphinxAtStartPar
comparison\_result\_o
\sphinxbeforeendvarwidth
\end{varwidth}%
&\begin{varwidth}[t]{\sphinxcolwidth{1}{3}}
\sphinxAtStartPar
1 bit
\sphinxbeforeendvarwidth
\end{varwidth}%
&\begin{varwidth}[t]{\sphinxcolwidth{1}{3}}
\sphinxAtStartPar
Comparator output
\sphinxbeforeendvarwidth
\end{varwidth}%
\\
\sphinxhline\begin{varwidth}[t]{\sphinxcolwidth{1}{3}}
\sphinxAtStartPar
is\_equal\_result\_o
\sphinxbeforeendvarwidth
\end{varwidth}%
&\begin{varwidth}[t]{\sphinxcolwidth{1}{3}}
\sphinxAtStartPar
1 bit
\sphinxbeforeendvarwidth
\end{varwidth}%
&\begin{varwidth}[t]{\sphinxcolwidth{1}{3}}
\sphinxAtStartPar
Equality flag
\sphinxbeforeendvarwidth
\end{varwidth}%
\\
\sphinxbottomrule
\end{tabulary}
\sphinxtableafterendhook\par
\sphinxattableend\end{savenotes}

\sphinxstepscope


\chapter{Functional Blocks}
\label{\detokenize{alu_functionalblocks:functional-blocks}}\label{\detokenize{alu_functionalblocks::doc}}
\sphinxAtStartPar
The ALU consists of four major blocks.


\section{Shared Adder}
\label{\detokenize{alu_functionalblocks:shared-adder}}
\sphinxAtStartPar
The shared adder performs:
\begin{itemize}
\item {} 
\sphinxAtStartPar
Addition

\item {} 
\sphinxAtStartPar
Subtraction

\item {} 
\sphinxAtStartPar
Signed comparisons

\item {} 
\sphinxAtStartPar
Unsigned comparisons

\end{itemize}

\sphinxAtStartPar
Instead of using separate hardware for addition and subtraction,
the ALU uses a single adder.

\sphinxAtStartPar
Subtraction is performed using two’s complement arithmetic:

\begin{sphinxVerbatim}[commandchars=\\\{\}]
\PYG{n}{A} \PYG{o}{\PYGZhy{}} \PYG{n}{B} \PYG{o}{=} \PYG{n}{A} \PYG{o}{+} \PYG{p}{(}\PYG{o}{\PYGZti{}}\PYG{n}{B} \PYG{o}{+} \PYG{l+m+mi}{1}\PYG{p}{)}
\end{sphinxVerbatim}


\section{Shifter}
\label{\detokenize{alu_functionalblocks:shifter}}
\sphinxAtStartPar
The shifter performs:
\begin{itemize}
\item {} 
\sphinxAtStartPar
Logical Left Shift (SLL)

\item {} 
\sphinxAtStartPar
Logical Right Shift (SRL)

\item {} 
\sphinxAtStartPar
Arithmetic Right Shift (SRA)

\end{itemize}

\sphinxAtStartPar
For left shifts, operand bits are first reversed and then shifted.


\section{Comparator}
\label{\detokenize{alu_functionalblocks:comparator}}
\sphinxAtStartPar
The comparator supports:

\sphinxAtStartPar
Equality Operations
\begin{itemize}
\item {} 
\sphinxAtStartPar
ALU\_EQ

\item {} 
\sphinxAtStartPar
ALU\_NE

\end{itemize}

\sphinxAtStartPar
Unsigned Comparisons
\begin{itemize}
\item {} 
\sphinxAtStartPar
ALU\_GTU

\item {} 
\sphinxAtStartPar
ALU\_GEU

\item {} 
\sphinxAtStartPar
ALU\_LTU

\item {} 
\sphinxAtStartPar
ALU\_LEU

\end{itemize}

\sphinxAtStartPar
Signed Comparisons
\begin{itemize}
\item {} 
\sphinxAtStartPar
ALU\_GTS

\item {} 
\sphinxAtStartPar
ALU\_GES

\item {} 
\sphinxAtStartPar
ALU\_LTS

\item {} 
\sphinxAtStartPar
ALU\_LES

\end{itemize}


\section{Result Multiplexer}
\label{\detokenize{alu_functionalblocks:result-multiplexer}}
\sphinxAtStartPar
The final output is selected according to operator\_i.

\sphinxAtStartPar
Supported operations include:
\begin{itemize}
\item {} 
\sphinxAtStartPar
ADD

\item {} 
\sphinxAtStartPar
SUB

\item {} 
\sphinxAtStartPar
AND

\item {} 
\sphinxAtStartPar
OR

\item {} 
\sphinxAtStartPar
XOR

\item {} 
\sphinxAtStartPar
SLL

\item {} 
\sphinxAtStartPar
SRL

\item {} 
\sphinxAtStartPar
SRA

\item {} 
\sphinxAtStartPar
EQ

\item {} 
\sphinxAtStartPar
NE

\item {} 
\sphinxAtStartPar
GT

\item {} 
\sphinxAtStartPar
GE

\item {} 
\sphinxAtStartPar
LT

\item {} 
\sphinxAtStartPar
LE

\item {} 
\sphinxAtStartPar
SLT

\item {} 
\sphinxAtStartPar
SLE

\end{itemize}

\sphinxstepscope


\chapter{ALU Architecture}
\label{\detokenize{alu_architecture:alu-architecture}}\label{\detokenize{alu_architecture::doc}}

\section{Internal Blocks}
\label{\detokenize{alu_architecture:internal-blocks}}\begin{enumerate}
\sphinxsetlistlabels{\arabic}{enumi}{enumii}{}{.}%
\item {} 
\sphinxAtStartPar
Shared Adder

\item {} 
\sphinxAtStartPar
Comparator

\item {} 
\sphinxAtStartPar
Shifter

\item {} 
\sphinxAtStartPar
Result Multiplexer

\end{enumerate}


\section{Architecture Diagram}
\label{\detokenize{alu_architecture:architecture-diagram}}
\sphinxstepscope


\chapter{Supported Operations}
\label{\detokenize{alu_operations:supported-operations}}\label{\detokenize{alu_operations::doc}}

\section{Arithmetic}
\label{\detokenize{alu_operations:arithmetic}}\begin{itemize}
\item {} 
\sphinxAtStartPar
\sphinxstylestrong{ADD} \sphinxhyphen{} Adds two operands together.

\sphinxAtStartPar
Example: 10 + 5 = 15

\item {} 
\sphinxAtStartPar
\sphinxstylestrong{SUB} \sphinxhyphen{} Subtracts the second operand from the first.

\sphinxAtStartPar
Example: 10 \sphinxhyphen{} 5 = 5

\end{itemize}


\section{Logical}
\label{\detokenize{alu_operations:logical}}\begin{itemize}
\item {} 
\sphinxAtStartPar
\sphinxstylestrong{AND} \sphinxhyphen{} Performs bitwise AND operation.

\sphinxAtStartPar
Example: 1100 AND 1010 = 1000

\item {} 
\sphinxAtStartPar
\sphinxstylestrong{OR} \sphinxhyphen{} Performs bitwise OR operation.

\sphinxAtStartPar
Example: 1100 OR 1010 = 1110

\item {} 
\sphinxAtStartPar
\sphinxstylestrong{XOR} \sphinxhyphen{} Performs bitwise Exclusive OR operation.

\sphinxAtStartPar
Example: 1100 XOR 1010 = 0110

\end{itemize}


\section{Shift}
\label{\detokenize{alu_operations:shift}}\begin{itemize}
\item {} 
\sphinxAtStartPar
\sphinxstylestrong{SLL} \sphinxhyphen{} Shifts bits to the left logically.

\sphinxAtStartPar
Example: 0001 << 2 = 0100

\item {} 
\sphinxAtStartPar
\sphinxstylestrong{SRL} \sphinxhyphen{} Shifts bits to the right logically and fills zeros.

\sphinxAtStartPar
Example: 1000 >> 2 = 0010

\item {} 
\sphinxAtStartPar
\sphinxstylestrong{SRA} \sphinxhyphen{} Shifts bits to the right while preserving the sign bit.

\sphinxAtStartPar
Example: 11111000 >> 2 = 11111110

\end{itemize}


\section{Comparison}
\label{\detokenize{alu_operations:comparison}}\begin{itemize}
\item {} 
\sphinxAtStartPar
\sphinxstylestrong{EQ} \sphinxhyphen{} Checks if two operands are equal.

\sphinxAtStartPar
Example: 5 == 5 → True

\item {} 
\sphinxAtStartPar
\sphinxstylestrong{NE} \sphinxhyphen{} Checks if two operands are not equal.

\sphinxAtStartPar
Example: 5 != 3 → True

\item {} 
\sphinxAtStartPar
\sphinxstylestrong{GT} \sphinxhyphen{} Checks if the first operand is greater than the second.

\sphinxAtStartPar
Example: 8 > 5 → True

\item {} 
\sphinxAtStartPar
\sphinxstylestrong{GE} \sphinxhyphen{} Checks if the first operand is greater than or equal to the second.

\sphinxAtStartPar
Example: 5 >= 5 → True

\item {} 
\sphinxAtStartPar
\sphinxstylestrong{LT} \sphinxhyphen{} Checks if the first operand is less than the second.

\sphinxAtStartPar
Example: 3 < 5 → True

\item {} 
\sphinxAtStartPar
\sphinxstylestrong{LE} \sphinxhyphen{} Checks if the first operand is less than or equal to the second.

\sphinxAtStartPar
Example: 5 <= 5 → True

\item {} 
\sphinxAtStartPar
\sphinxstylestrong{SLT} \sphinxhyphen{} Sets result if the first signed operand is less than the second.

\sphinxAtStartPar
Example: \sphinxhyphen{}2 < 3 → True

\item {} 
\sphinxAtStartPar
\sphinxstylestrong{SLE} \sphinxhyphen{} Sets result if the first signed operand is less than or equal to the second.

\sphinxAtStartPar
Example: \sphinxhyphen{}2 <= \sphinxhyphen{}2 → True

\end{itemize}

\sphinxstepscope


\chapter{Design Features}
\label{\detokenize{alu_designfeatures:design-features}}\label{\detokenize{alu_designfeatures::doc}}\begin{itemize}
\item {} 
\sphinxAtStartPar
Shared hardware architecture

\item {} 
\sphinxAtStartPar
Fast execution

\item {} 
\sphinxAtStartPar
Reduced area

\end{itemize}

\sphinxstepscope


\chapter{Conclusion}
\label{\detokenize{alu_conclusions:conclusion}}\label{\detokenize{alu_conclusions::doc}}
\sphinxAtStartPar
The zeroriscy\_alu module is the execution unit of the
Zero\sphinxhyphen{}RISCY processor.

\sphinxAtStartPar
It efficiently performs arithmetic, logical,
comparison and shift operations using shared hardware
resources.



\renewcommand{\indexname}{Index}
\printindex
\end{document}